\section{Heat-pump performance}
\label{app:cop}

The heat-pump COPs are calculated exogenously with an adaptation of the central-heating COP approximation implemented in PyPSA-Eur \cite{pypsaeur2025}, which is based on Jensen et al.\ \cite{jensen2018cop} with the parameterisation evaluated by Pieper et al.\ \cite{pieper2020cop}.

\paragraph{Supply temperature}
The district-heating supply temperature follows HTF's weather-compensation curves. For each of three network points $n$, the required temperature decreases linearly with the 24-hour mean ambient temperature $\overline{T}^{\mathrm{amb}}_t$ down to a minimum value,
\begin{equation}
    T^{\mathrm{sup}}_{n,t} = \max\left\{T^{\min}_n,\ T^{0}_n - \frac{T^{0}_n - T^{\min}_n}{16}\,\overline{T}^{\mathrm{amb}}_t\right\},
\end{equation}
with $(T^{\min}_n, T^{0}_n)$ equal to $(66, 74)$, $(68, 78)$ and $(63, 73)$~$^\circ$C. The supply temperature is the highest of the three, $T^{\mathrm{sup}}_t=\max_n T^{\mathrm{sup}}_{n,t}$, and ranged from 68 to 83~$^\circ$C in 2024. The return temperature is fixed at 50~$^\circ$C. Ambient temperature is taken from ERA5-Land.

\paragraph{COP}
For each heat pump $i\in\{\mathrm{c2},\mathrm{gw}\}$, the logarithmic mean temperatures of sink and source are
\begin{equation}
    \widetilde{T}^{i}_{j,t}
    = \frac{T^{i,\mathrm{hot}}_{j,t} - T^{i,\mathrm{cold}}_{j,t}}
           {\ln\left(T^{i,\mathrm{hot}}_{j,t}/T^{i,\mathrm{cold}}_{j,t}\right)},
    \qquad j\in\{\mathrm{sink},\mathrm{src}\},
\end{equation}
and the ideal Lorenz COP is
\begin{equation}
    COP^{\mathrm{L},i}_{t} = \frac{\widetilde{T}^{i}_{\mathrm{sink},t}}{\widetilde{T}^{i}_{\mathrm{sink},t} - \widetilde{T}^{i}_{\mathrm{src},t}}.
\end{equation}
It is adjusted for refrigerant properties, the heat-exchanger pinch temperature and compressor efficiency:
\begin{equation}
    COP^{i}_{t} = COP^{\mathrm{L},i}_{t} F^{i}_{t}\,\eta^{\mathrm{is}}\left(1-r^{i}_{t}\right) + \left(1-\eta^{\mathrm{is}}\right) - l^{\mathrm{HP}},
\end{equation}
where $F^{i}_{t}$ and $r^{i}_{t}$ are the correction factors of the approximation. Ammonia is assumed as refrigerant, with an isentropic compressor efficiency of 0.75, a pinch-point temperature difference of 2~K and no additional heat-loss term. COP values below one or temperature lifts below 5~K are excluded \cite{zajacs2020lowtemperature}.

The source inlet and outlet temperatures are 12/8~$^\circ$C for City2 and 10.5/4~$^\circ$C for the groundwater heat pump. The resulting hourly COPs range from 3.14 to 3.56 for City2 and from 2.97 to 3.34 for the groundwater heat pump.
